\documentclass[11pt]{article}
\usepackage[a4paper,margin=25mm]{geometry}
\usepackage[T1]{fontenc}
\usepackage[utf8]{inputenc}
\usepackage{lmodern}
\usepackage{microtype}
\usepackage{amsmath,amssymb}
\usepackage{graphicx}
\usepackage{booktabs}

\usepackage[numbers,sort&compress]{natbib}
\usepackage[colorlinks=true,linkcolor=black,citecolor=black,urlcolor=blue]{hyperref}
\usepackage{caption}

\captionsetup{font=small,labelfont=bf,width=0.95\textwidth}
\graphicspath{{../outputs/paper/tex/}}
\newcommand{\code}[1]{\texttt{#1}}
\newcommand{\Rsq}{R^{2}}

\title{What the correlations between neurons say about each neuron}
\author{Francesco L\"assig\thanks{TODO: affiliation and co-authors.}}
\date{Draft v2, \today}

\begin{document}
\maketitle

\begin{abstract}
Record many neurons while they watch the same stimulus, then discard the stimulus. What remains is how the neurons relate to each other: a correlation matrix. We ask whether that matrix alone says what each neuron represents. From in-vivo responses in the MICrONS volume, a transformer that sees only the correlation matrix of 512 sampled neurons, with no labels and no reference population, decodes the preferred orientation of mouse visual cortex neurons to a mean error of $25^\circ$ (uninformed: $45^\circ$; best constant guess given the label distribution: $41^\circ$; $28^\circ$ when every orientation is weighted equally) and their receptive-field position, scored against twin-model estimates, at $\Rsq = 0.23$. On the digital-twin substrate it reaches $20^\circ$ ($24^\circ$ balanced) and 0.48. Asked whether a neuron prefers horizontal or vertical, the decoder is right for 79\% of neurons in vivo and 87\% on the twin, against 56\% for the majority answer. It does not resolve the obliques. The readout is possible because the relations between orientation classes are close to circulant, which fixes the class structure up to rotation and reflection. The absolute frame comes from a departure from that symmetry. The departure survives equal representation of every orientation in each population and disjoint training and test stimuli, and it coincides with an anisotropy of the correlations: in each area the neurons preferring the dominant orientation co-fluctuate most strongly, and a model population with only that property, at the strength measured in cortex, recovers the frame. Across 33 mice of the Allen Brain Observatory, with relations from drifting gratings, orientation transfers to unseen animals only when a sampled population mixes animals, so that its matrix contains between-animal correlations. Populations drawn from one animal carry only a weak signal, $2^\circ$ to $3^\circ$ below chance. That reflects the training diversity one animal provides, not its circuit: single-scan populations in MICrONS are decoded to $31^\circ$ to $33^\circ$. For receptive fields, what crosses animals is where each animal's imaged patch looks on the screen, not the layout of neurons within it.
\end{abstract}

\section{Introduction}

What a neuron in visual cortex represents is established by relating its activity to the stimulus: a neuron prefers vertical gratings, or responds to the upper left of the screen. We ask whether the same facts are contained in a description that mentions no stimulus. Take a population that watched the same movie, compute the correlation of every pair of neurons, and discard the movie. Does the resulting matrix determine which neuron prefers which orientation, and where each neuron's receptive field lies?

The idea that content can be read from relations is old. Shepard and Chipman proposed that internal representations need not resemble what they stand for as long as the similarities among them mirror the similarities among the things represented, a second-order isomorphism \citep{shepard1970}. Edelman argued that representation is representation of similarities \citep{edelman1998}. Representational similarity analysis made it a method that compares brains, models and species through their similarity structures rather than their coordinates \citep{kriegeskorte2008,kriegeskorte2013}, hyperalignment builds a shared representational space across subjects from those structures \citep{haxby2011}, and the same principle underlies current work on aligning representations across networks and individuals \citep{sucholutsky2023,huh2024}. In neuroscience, the correlation structure of spontaneous activity was shown to carry the orientation map of visual cortex without any stimulus \citep{tsodyks1999,kenet2003} and, more generally, to match the statistics of evoked activity \citep{berkes2011}; functionally similar neurons connect preferentially \citep{ko2011,ding2025}; and population dynamics are preserved across animals performing the same behaviour \citep{safaie2023}. Structuralist accounts of experience claim that what a state represents, or how it feels, is fixed by its position in a space of relations \citep{lyre2022}, with \citet{kleiner2024} giving a formal definition of such structures. \citet{framework2025} formulates the corresponding question for neural networks and gives the decoding framework we use here.

The closest precedent is the spontaneous-activity work: a neuron's correlations with the population contain its tuning even in the dark. \citet{tsodyks1999} showed it for single neurons by averaging the imaged population activity at the times a neuron fired and finding the evoked map of its preferred orientation. Reading the tuning out of that pattern required the evoked map, a stimulus-labelled reference. Correlations sort neurons by orientation and place the groups in order around the orientation circle, but on their own they say nothing about which group is vertical. The same holds for representational similarity analysis and hyperalignment, whose structures are over conditions or subjects and whose interpretation goes through a known stimulus set. We ask instead whether the relations among \emph{neurons}, with the stimulus discarded and no labelled neuron or map to compare with, fix the content of each neuron, and whether that content is readable across thousands of individually identified neurons and across animals. The nearest machine-learning work on neural data reads properties of a neuron from population activity: NeuPRINT learns a time-invariant embedding of each neuron from the dynamics of its population and decodes cell type from it \citep{mi2023}, NuCLR does the same with a contrastive objective and a permutation-equivariant transformer, and transfers cell type and brain region to unseen animals \citep{arora2025}, and POYO trains one decoder across sessions and animals by tokenising spikes with learned unit embeddings \citep{azabou2023}. These models see the activity itself. Our decoder sees only the correlation matrix, so it answers what the relations alone determine, and its target is stimulus tuning rather than cell class. The relational situation is often the one at hand: recordings under natural behaviour have no controlled stimulus, and neurons from different sessions or animals share no stimulus-aligned coordinates, but they do share the structure of their correlations. The answer also says how much of a neuron's content is fixed by its position in the relations among neurons.

We answer it with a decoder. The input is the correlation matrix of a sampled population of a few hundred neurons, one row per neuron, and nothing else: no labels and no reference neurons. The output is a preferred orientation and a receptive-field position for every neuron in the population. The decoder is trained on populations whose labels are known and scored on neurons it never saw. We do this at two scales. Within one brain, using the MICrONS release \citep{ding2025,microns2025}, 12,894 neurons from a cubic millimetre of mouse visual cortex with in-vivo responses to a shared movie, a fitted digital-twin model \citep{wang2025}, and labels for orientation and receptive field (Fig.~\ref{fig:pipeline}): the decoder trained on half the neurons reads the other half. Across brains, using the Allen Brain Observatory \citep{devries2020}, 33 mice recorded under identical stimuli: the decoder trained on some animals reads neurons of animals it never saw.

Averaging the correlations by orientation class shows that the relations between classes are almost circulant, which fixes the classes only up to a rotation or reflection of the orientation circle. A small, consistent departure from that symmetry fixes the absolute frame and makes the readout possible.

\begin{figure}[tbp]
  \centering
  \includegraphics[width=\textwidth]{fig1_microns_pipeline.pdf}
  \caption{\textbf{How the MICrONS task is built.} (a)~One mouse, 13 two-photon scans of a cubic millimetre of visual cortex (areas V1, RL, AL, LM), 12,894 neurons co-registered to the electron-microscopy volume. (b)~The relational substrate. Each neuron is its trial-averaged response vector to a stimulus all neurons saw (in vivo: 120 bins of the oracle natural-movie clips; digital twin: 4,999 bins compressed to 512 principal components); shown are six real neurons and their correlation matrix. The oracle clips were shown in all 13 scans, so the matrix is defined across scans. The decoder sees only such correlations, never the stimulus. (c)~The contents of the same six neurons: preferred orientation (in vivo, a continuous angle, 5,287 neurons with gOSI $\geq 0.25$) and receptive-field centre (digital-twin fit, 11,326 neurons). These are training targets only. (d)~The protocol, with two terms kept apart. The \emph{halves}: all neurons are split at random into a training half and a test half, and the full $12{,}894 \times 12{,}894$ correlation matrix, ordered training-first (a real 200-neuron excerpt is shown), is used only in its training$\times$training and test$\times$test blocks. A \emph{sample}: 512 neurons drawn at random from one half; its Gram is the corresponding diagonal sub-block (dashed), and that $512\times512$ matrix is the decoder's entire input, one row per neuron, one prediction per row. Training draws thousands of samples per epoch from the training half; a loss compares their predictions with the labels of the labelled neurons in the sample. Scoring draws samples from the test half and compares predictions with held-out labels; no gradient. Labels never enter the input on either path, and the halves share no neuron. The epoch is chosen on a held-out slice of the training half.}
  \label{fig:pipeline}
\end{figure}

\section{Methods}

\subsection{Data}

\paragraph{MICrONS.} The release of \citet{ding2025} contains 12,894 neurons co-registered to the electron-microscopy volume, from 13 two-photon scans across V1, RL, AL and LM. Each neuron has a trial-averaged in-vivo response to the oracle natural-movie clips (120 bins), shown in every scan; a digital-twin response to a shared movie (4,999 bins, which we compress to 512 principal components; the Gram is preserved to $r > 0.999$); an in-vivo preferred orientation with a global orientation selectivity index (gOSI); and a receptive-field centre from the twin (spike-triggered-average fit, screen coordinates in $[-1,1]$). We use orientation, as a continuous angle, for the 5,287 neurons with gOSI $\geq 0.25$, and receptive-field position for the 11,326 neurons whose twin test correlation is $\geq 0.2$. The orientation labels are in-vivo measurements on both substrates; the receptive-field labels come from the twin's fit on both substrates, so every receptive-field result, in vivo included, scores relations against model-derived positions (Section~4). Same-scan pairs correlate about 50\% more strongly than cross-scan pairs at matched receptive-field distance. The content of the correlations is the same across scans, and all populations mix scans. Only 2\% of the receptive-field label variance lies between scans and 9\% between areas, so a decoder that read scan or area identity from the Gram could reach at most that much $\Rsq$.

\paragraph{Allen.} From Visual Coding 2P \citep{devries2020} we take the 36 VISp excitatory containers with at least 150 orientation-labelled cells: 33 mice, 9,281 cells, of which 5,529 carry an orientation label; unlabelled cells stay in the populations and receive no loss, so a held-out-animal split has 6,706 training cells and 2,575 test cells in total. Relations come from session A, either responses to natural movies shown to every mouse (4,500 bins) or the 40 blank-subtracted condition means of drifting gratings (8 directions $\times$ 5 temporal frequencies). Orientation labels come from static gratings in session B (six orientations $30^\circ$ apart, gOSI $\geq 0.25$, $p < 0.05$), so labels and relations use different stimuli. Receptive-field centres come from locally sparse noise in session C (1,504 cells with a significant fit, 23 to 112 per mouse).

\subsection{Relations}

Each neuron's response vector is centred and normalised; the relation between two neurons is the cosine of their vectors, a signal correlation. For a population of $n$ neurons the Gram is the $n\times n$ matrix of these values, standardised by the global off-diagonal mean and standard deviation, with the diagonal zeroed. We call this standardised correlation matrix the Gram throughout.

\subsection{Decoder}

Each row of a population's Gram, together with three row statistics (the mean, standard deviation and third moment of the row), is projected to a token; $L$ pre-norm transformer blocks \citep{vaswani2017} with $H$-head attention update the tokens, the Gram entering the attention logits as an additive bias; a linear head reads out, per token, the unit vector $(\cos 2\theta, \sin 2\theta)$ of the preferred orientation $\theta$ or an $(x,y)$ receptive-field position. The loss is the mean squared error on the labelled tokens only; unlabelled neurons stay in the population and receive no loss. The standard configuration is width 512, 8 layers and 17M parameters with populations of 512 neurons on MICrONS, and width 256, 4 layers and 2.2M parameters with populations of 256 neurons on Allen. Other sizes are in the appendix. Training draws thousands of random populations per epoch.

\subsection{Protocol}

\paragraph{Split.} MICrONS neurons are halved at random. Training populations are sampled inside the training half and test populations inside the test half, so no test neuron appears in a training population (Fig.~\ref{fig:pipeline}d). A slice of the training half is held out to choose the epoch. The test half never influences selection.

\paragraph{Regimes (Allen).} Two independent choices define a run (Fig.~\ref{fig:regimes}). The split decides where test neurons come from: the training animals, each halved into training and test neurons, or held-out animals. The population decides what a sample is: neurons drawn from one animal, so that the Gram is a within-circuit matrix, or from several animals, so that most entries are between-animal correlations. We call the four cells \code{within}, \code{pooledwithin}, \code{cross} and \code{pooledcross}. Held-out regimes select the epoch on further held-out mice. Every held-out cell is run on three mouse splits and every training-animal cell on three neuron splits.

\paragraph{Metrics.} Orientation is scored by the mean absolute angular error between the decoded angle, $\tfrac12\operatorname{atan2}$ of the output, and the label, taken modulo $180^\circ$, so it lies between $0^\circ$ and $90^\circ$. A decoder whose output is independent of the label produces errors spread evenly over that range when either the labels or the outputs are uniform on the orientation circle, so chance is $45^\circ$. With the MICrONS label distribution the best constant prediction scores $41.3^\circ$ (Appendix~\ref{app:controls}). Both lines are drawn in the figures. Allen labels sit on a $30^\circ$ grid, so a decoder that recovered every true preference exactly would still disagree with them by $7.5^\circ$ on average. That is the floor for Allen errors. A frame check reports the error after the best single rotation or reflection of all decoded angles. If it is lower than the raw error, the decoder recovered the structure but not the frame. Receptive fields are scored by $\Rsq$ against a baseline of 0.

\subsection{Symmetry analysis}

To expose the symmetry of the relations we average the Gram into a $K\times K$ class-Gram, the mean relation between neurons of two orientation classes, project it onto its circulant part, and ask which relabellings a held-out half's class-Gram is compatible with (all $K!$ relabellings, Frobenius distance). A relabelling that differs from the truth by a rotation or a reflection of the orientation circle is an element of the dihedral group $D_K$. The analysis reports whether the truth is singled out absolutely or only modulo $D_K$.

\section{Results}

\subsection{Per-neuron content from the correlation matrix alone}

On MICrONS the decoder reads orientation and receptive-field position from the Gram of 512 neurons of the test half (Fig.~\ref{fig:decoder}, Table~\ref{tab:microns}). The twin substrate is cleaner, since the twin removes trial noise, and receptive fields carry more than orientation. The predictions are in the true frame: the frame check lowers the error by less than $0.1^\circ$ in every run. About half the neurons are decoded to within $15^\circ$ (0.46 in vivo, 0.54 on the twin), but the error is not spread evenly over the orientation circle: the decoder resolves the cardinal axes and little else (Appendix~\ref{app:controls}, Fig.~\ref{fig:bins}). Asked only which cardinal axis a neuron is closer to, it is right for 87\% of all neurons on the twin and 79\% in vivo (majority answer 56\%; averaged over the eight true-orientation bins, 83\% and 74\%), and for the 56\% of neurons whose preference lies within $15^\circ$ of an axis, for 94\% and 86\% (neuron splits 1 and 2). Sorted into four classes at $0^\circ$, $45^\circ$, $90^\circ$ and $135^\circ$, 66\% (twin) and 59\% (in vivo) land in the right one against a majority class of 40\%. The misses are almost all obliques. Within $15^\circ$ of a cardinal axis the error is $12^\circ$ on the twin and $19^\circ$ in vivo; within $15^\circ$ of an oblique, $36^\circ$ to $39^\circ$, near chance (Fig.~\ref{fig:bins}). Both readouts are properties of individual neurons, not of the scan: centring predictions and labels per scan leaves the receptive-field $\Rsq$ unchanged and the orientation error within $0.2^\circ$ (Appendix~\ref{app:controls}). Test Grams from stimulus bins the decoder never trained on add about $0.5^\circ$ of error on the twin and $8.5^\circ$ in vivo, where a Gram over 60 movie bins is partly specific to those bins (Appendix~\ref{app:controls}).

\begin{table}[tbp]
  \centering\footnotesize
  \caption{MICrONS, per-neuron decoding on the test half. Relations only: the 17M decoder on the Gram of 512-neuron populations, mean $\pm$ s.d.\ over 3 seeds. Balanced: the error averaged over eight equally weighted true-orientation bins (neuron split 2). Stimulus-known linear readout: a linear map from each neuron's own movie response vector to its content, a reference that uses the stimulus alignment the relational decoder discards, not a label-noise ceiling (Appendix~\ref{app:controls}). Chance: an uninformed decoder ($45^\circ$) and the best constant orientation given the label distribution ($41.3^\circ$).}
  \label{tab:microns}
  \begin{tabular}{llcccc}
    \toprule
    content & substrate & relations only & \begin{tabular}[c]{@{}c@{}}balanced over\\orientations\end{tabular} & \begin{tabular}[c]{@{}c@{}}stimulus-known\\linear readout\end{tabular} & \begin{tabular}[c]{@{}c@{}}chance /\\best constant\end{tabular} \\
    \midrule
    orientation (angular error) & in vivo & $25.4^\circ \pm 0.5^\circ$ & $27.9^\circ$ & $20.2^\circ$ & $45^\circ$ / $41.3^\circ$ \\
    orientation (angular error) & twin & $20.2^\circ \pm 0.4^\circ$ & $24.5^\circ$ & $15.8^\circ$ & $45^\circ$ / $41.3^\circ$ \\
    receptive field ($\Rsq$) & in vivo & $0.234 \pm 0.007$ & -- & 0.34 & 0 \\
    receptive field ($\Rsq$) & twin & $0.480 \pm 0.010$ & -- & 0.68 & 0 \\
    \bottomrule
  \end{tabular}
\end{table}

\begin{figure}[tbp]
  \centering
  \includegraphics[width=0.85\textwidth]{fig2_microns_decoder.pdf}
  \caption{\textbf{MICrONS: per-neuron content decoded from the population correlation matrix alone.} Left: mean angular error of the decoded preferred orientation on held-out neurons of the same animal, for decoders of 0.3M, 2.2M and 17M parameters; error bars are the s.d.\ over 3 seeds; grey line: the $45^\circ$ error of an uninformed decoder; black marks: the stimulus-known linear readout of each neuron's own responses. Middle: $\Rsq$ of the decoded receptive-field centre, same marks. Right: coarse readouts of the same decoder on neuron splits 1 and 2: whether the decoded angle falls on the correct cardinal axis, for all neurons and for the neurons whose preference lies within $15^\circ$ of an axis, and whether it falls in the correct one of four classes at $0^\circ$, $45^\circ$, $90^\circ$ and $135^\circ$; black ticks: majority class. The decoder sees only the standardised correlation matrix of 512 sampled neurons.}
  \label{fig:decoder}
\end{figure}

\begin{figure}[tbp]
  \centering
  \includegraphics[width=0.95\textwidth]{figA5_error_by_orientation.pdf}
  \caption{\textbf{The readout is cardinal.} Mean angular error of the neurons in each true-orientation bin (bars; grey line: $45^\circ$ chance; dotted: best constant), with the fraction of labels (black) and of predictions (red) per bin, for the twin and in-vivo MICrONS decoders (neuron split 2) and the Allen \code{pooledcross} decoder. Cardinal orientations are decoded well, obliques near chance, and most predictions fall on $0^\circ$ or $90^\circ$.}
  \label{fig:bins}
\end{figure}

\subsection{Why the absolute frame is recoverable}

The class-Gram of orientation is close to circulant (Fig.~\ref{fig:symmetry}): neighbouring orientations correlate, orthogonal ones anti-correlate, and the pattern repeats around the circle. A circulant matrix is unchanged by every rotation and reflection of the orientation circle, so that part of the relations fixes the classes only up to the 16 elements of $D_8$ and cannot say which class is $0^\circ$. Whatever fixes the frame must be a departure from circulant structure, at the class level or finer. At the class level that residual holds 29\% of the class-Gram variance in vivo and 19\% on the twin, and it is signal: it correlates at 0.94 between random halves of the neurons on both substrates, at 0.84 (in vivo) and 0.99 (twin) between disjoint halves of the stimulus bins, and at 0.99 with the residual computed on class-balanced subsamples, so it is neither a finite-stimulus artefact nor a consequence of the uneven number of neurons per orientation. Its form is a cardinal anisotropy, stronger relations among neurons near the cardinal axes, whose shape differs between areas (residual correlation V1--RL 0.77, V1--AL 0.31).

Five observations locate the structure in the circulant part and narrow down what fixes the frame. The last two point to a specific anisotropy, which a model population then tests.

\emph{Priors do not explain the frame.} The decoder recovers the absolute frame (Section~3.1), while predictors that use only the distribution of labels, the best constant orientation or each scan's or area's mean, reach $39^\circ$ to $41^\circ$.

\emph{The structure is recoverable with no frame in the labels.} A decoder trained with each population's targets rotated at random, and the loss taken after the best alignment of predictions to targets, places neurons to $19.9^\circ$ on the twin and $25.5^\circ$ in vivo after one global alignment, the same as the standard decoder with the frame ($20.2^\circ$, $25.4^\circ$). The per-neuron structure needs no frame anywhere in the supervision. This does not show that the circulant part alone carries the structure, since the residual is still in that decoder's input. Its raw error, $25^\circ$ and $31^\circ$ against $45^\circ$ for an arbitrary frame, shows that it fixes a consistent frame from the data even when the loss does not ask for one.

\emph{The frame is not the label density and not the stimulus sample.} With every training and test population drawn with equal numbers of neurons from each of the eight orientation classes, the decoder reaches $22.0^\circ$ on the twin and $26.6^\circ$ in vivo (two seeds each), with no gap between raw and aligned error. It is better than the standard decoder scored with the same equal weighting ($24.5^\circ$, $27.9^\circ$). The frame also survives training and test relations built from disjoint halves of the stimulus (Appendix~\ref{app:controls}).

\emph{The errors have the shape of the symmetry breaking.} They are $11^\circ$ to $12^\circ$ near the cardinal axes and $37^\circ$ to $40^\circ$ at the obliques on the twin, and reweighting the loss so that obliques count as much as cardinals leaves them unchanged (Appendix~\ref{app:controls}). The oblique ordering is not in the relations at single-neuron resolution.

\emph{Structure transfers between areas; the frame does not reliably.} A decoder trained on V1 transfers to RL, whose anisotropy resembles V1's, with a $1.6^\circ$ frame gap. On AL, whose anisotropy differs, its raw error of $58^\circ$ falls to $32^\circ$ after a global rotation of $78^\circ$. That rotation is not evidence about AL's correlations: the V1 decoder predicts within $15^\circ$ of $90^\circ$ for every AL neuron, so the rotation only moves a pile of identical predictions onto AL's dominant orientation, $0^\circ$. With class-balanced training populations, which remove the prior, the predictions spread, transfer to AL is weak ($38^\circ$ and $53^\circ$ raw, $34^\circ$ and $37^\circ$ aligned in two seeds), and the rotation the two seeds need differs ($0^\circ$ and $85^\circ$); RL keeps V1's frame ($30^\circ$ raw and aligned). The class-level residuals of V1 and AL do not single out any $D_8$ element either (best correlation 0.31 on the twin, flat across elements; AL has 294 labelled neurons). %V3PENDING-BT
\emph{The dominant orientation of an area is its most strongly co-correlated one.} Define the anisotropy factor as the mean correlation between two neurons of the same orientation class, at the class where it peaks, divided by its mean over the other seven classes. It peaks at $90^\circ$ in V1 on both substrates (factor 1.28 on the twin, 0.071 against 0.055; 1.77 in vivo, 0.095 against 0.054), at $0^\circ$ in AL on the twin (1.49) and at $135^\circ$ in AL in vivo (1.69), and the values are unchanged on class-balanced subsamples, so the signature is in the correlations, not in the counts. The three row statistics alone, which include each neuron's mean correlation, fix a weak frame by themselves ($37.7^\circ$ raw and aligned; Appendix~\ref{app:controls}). Cardinal over-representation and its variation between areas are established features of mouse visual cortex \citep{roth2012,fahey2019}, and the co-fluctuation of like-tuned neurons follows from their preferential connectivity \citep{ko2011} and from the stimulus dependence of noise correlations \citep{kohn2005}.

\emph{A model population tests the candidate.} We could not isolate the anchor by subtraction, because removing the class-level residual from every Gram with the true classes injects the classes into the Gram (Appendix~\ref{app:controls}), so the frame is located by exclusion, by the cross-area pattern, and by a model (Appendix~\ref{app:controls}, Table~\ref{tab:syn}). Model neurons with von Mises tuning, a movie-like stimulus that covers orientation evenly, and the same pipeline give a class-Gram circulant to within noise, and the decoder trained on it recovers the ordering, $30^\circ$ after one global alignment, but no frame, $54^\circ$ raw. Sharper cardinal tuning and a cardinal excess of neurons fix the frame in at most one seed of five, alone, together, or at any strength tested; an excess of cardinal frames in the stimulus does in three of five. Adding only the anisotropy the data show, stronger co-fluctuation among like-tuned neurons peaked at one orientation, with uniform counts and an even stimulus, fixes the frame: at a factor of 1.32, the twin's, in four of five seeds at $9^\circ$ raw; at 1.75, in vivo V1's, in five of five at $3^\circ$; at 1.09, below any cortical value, in one of two. %V3PENDING-SYN
The cortical anisotropy is not inherited from the stimulus, since it reproduces across disjoint halves of the stimulus and differs between areas that saw the same movie.

\begin{figure}[tbp]
  \centering
  \includegraphics[width=0.9\textwidth]{fig3_symmetry.pdf}
  \caption{\textbf{What fixes the absolute orientation frame.} Left: mean correlation between orientation classes in vivo (diagonal masked). Middle: its circulant projection. Right: matching a held-out half's class-Gram to the reference over all $8!$ relabellings. With the raw matrix the true labelling wins. After the circulant projection, absolute accuracy collapses to chance while accuracy modulo $D_8$ (dots) stays at 1.0, since a circulant matrix is invariant under $D_8$. The bars sit at or slightly below the chance line because the matching then picks one of the 16 equivalent labellings at random. The panel shows what the residual has to supply.}
  \label{fig:symmetry}
\end{figure}

\subsection{Orientation transfers between animals through mixed populations}

With natural-movie relations, orientation is nearly absent from Allen correlations: same-orientation pairs correlate 0.005 more than orthogonal pairs, against 0.09 in MICrONS, and every test is at chance. The Allen results below therefore use grating relations, a designed set of 40 conditions over an orientation grid, and do not replicate on naturalistic relations. With gratings the signal is present (0.05 within a mouse, 0.044 across mice) and the four regimes separate (Fig.~\ref{fig:allen2x2}, Table~\ref{tab:allen}). Single-animal populations carry a weak but consistent signal, $1.9^\circ$ to $2.9^\circ$ below chance in all six runs, whether or not the animal was seen in training. Mixed populations are $6.6^\circ$ to $9.8^\circ$ below chance in both rows, at the same level, so no loss from transfer to a new brain is detectable at three splits. As on MICrONS the readout is cardinal: over the three splits the decoder names the right cardinal axis for 62\% of all cells and 68\% of the cells labelled $0^\circ$ or $90^\circ$ (majority 51\% and 52\%; 61\% averaged over the six label bins), cells labelled $0^\circ$ are decoded to $25^\circ$ and those labelled $90^\circ$ to $34^\circ$, and cells labelled $120^\circ$ or $150^\circ$ are at chance (Fig.~\ref{fig:bins}). What matters is whether the Gram contains between-animal correlations. These exist only because all mice saw the same 40 conditions, and they measure how similar two tuning profiles are. The decoder still never sees which condition is which, so the absolute frame must again come from the population's relational structure. Orientation is thus readable from a correlation geometry that spans animals, and only weakly from relations inside one circuit. Two controls say why. The limit is not the neurons: with each neuron's own 40-condition response vector as its token, the same decoder on the same single-animal populations reaches $33^\circ$ to $35^\circ$, as with mixed populations (Table~\ref{tab:allen}, last row). The limit is the training pool: mixed populations drawn from only 280 or 560 training cells, one or two animals' worth, score $40.5^\circ$ and $40.7^\circ$, near chance, against $36.0^\circ$ from the full pool of 6,706 at the same population size of 128 (Appendix~\ref{app:controls}). An animal has about 280 labelled cells, half of them training cells in the within-animal split, so a decoder restricted to within-animal correlations is trained on about 140 cells per animal, too few to learn from whatever those correlations contain. The same limit appears on MICrONS: a decoder trained on 140 neurons of one scan and tested on the other scans scores $38^\circ$ to $44^\circ$ on the twin and $37^\circ$ to $39^\circ$ in vivo (three scans; Appendix~\ref{app:controls}). Where one animal supplies thousands of training neurons, single-scan populations are decoded to $30.7^\circ$ on the twin and $32.6^\circ$ in vivo (128 neurons; Appendix~\ref{app:controls}), against $22^\circ$ for mixed-scan populations of the same size. Relations inside one circuit do carry orientation, more weakly than relations that span circuits.

\begin{table}[tbp]
  \centering\footnotesize
  \caption{Allen, mean angular error of the decoded orientation in the four regimes (2.2M decoder, 256-neuron populations; 128 for single-animal populations), three splits per cell (of neurons within each animal for the top row, of animals for the second row). Last row: the same decoder on held-out animals with each neuron's 40-condition response vector as its token instead of its Gram row, the stimulus-known reference. Chance $45^\circ$, best constant given the Allen label distribution $43.9^\circ$; label grid $7.5^\circ$. Means over the three splits in parentheses.}
  \label{tab:allen}
  \begin{tabular}{lcc}
    \toprule
    test neurons from & single-animal populations & mixed populations \\
    \midrule
    the training animals & $42.5^\circ$ / $42.9^\circ$ / $42.1^\circ$ ($42.5^\circ$) & $35.3^\circ$ / $37.2^\circ$ / $36.1^\circ$ ($36.2^\circ$) \\
    held-out animals & $43.1^\circ$ / $43.1^\circ$ / $42.7^\circ$ ($43.0^\circ$) & $35.2^\circ$ / $37.0^\circ$ / $38.4^\circ$ ($36.9^\circ$) \\
    \midrule
    held-out animals, stimulus-known reference & $32.8^\circ$ / $33.6^\circ$ / $34.9^\circ$ ($33.8^\circ$) & $32.4^\circ$ / $33.7^\circ$ / $34.3^\circ$ ($33.5^\circ$) \\
    \bottomrule
  \end{tabular}
\end{table}

\begin{figure}[tbp]
  \centering
  \includegraphics[width=0.95\textwidth]{fig4_regimes.pdf}
  \caption{\textbf{The four Allen decoder regimes.} Rows: test neurons from the training animals (each animal's cells halved) or from held-out animals. Columns: each sampled population drawn from one animal or from several. Filled dots are training neurons, hollow dots test neurons, dashed boxes one sampled population each. Training populations come only from training neurons and test populations only from test neurons in every regime.}
  \label{fig:regimes}
\end{figure}

\begin{figure}[tbp]
  \centering
  \includegraphics[width=0.85\textwidth]{fig5_allen_2x2_orientation.pdf}
  \caption{\textbf{Allen, 33 mice, grating relations: orientation is readable only from populations that mix animals.} Left: mean angular error; bars: mean over splits; points: the three splits; solid line: $45^\circ$ chance; dashed line: the $7.5^\circ$ label grid. Single-animal populations are $2^\circ$ to $3^\circ$ below chance whether the animal was seen in training or not; mixed populations are $7^\circ$ to $10^\circ$ below it in both rows. Right: coarse readouts of the mixed-population decoder on the held-out mice (three splits): correct cardinal axis for all cells and for the cells labelled $0^\circ$ or $90^\circ$; the third bar, nearest of four classes $45^\circ$ apart, maps each $30^\circ$ label to one class ($30^\circ$ and $60^\circ$ to $45^\circ$, $120^\circ$ and $150^\circ$ to $135^\circ$) and is shown for comparison with MICrONS; black ticks: majority class.}
  \label{fig:allen2x2}
\end{figure}

\subsection{Receptive fields: what crosses animals is the animal's screen position}

Correlation falls with receptive-field distance within a mouse (0.127 at $<5^\circ$ to 0.068 beyond $40^\circ$) and across mice (0.049 to 0.011). Decoding absolute screen coordinates on held-out mice gives $\Rsq$ that varies with the split, from $0.29$ to $-0.45$ for single-animal populations and from $0.21$ to $-0.25$ for mixed ones. Decomposing the saved predictions explains the variation (Fig.~\ref{fig:rf}). The correlation between predicted and true mean position of each held-out mouse, pooled over 27 mouse-level predictions, is 0.36 to 0.71 (permutation $p \leq 0.03$). The correlation between predicted and true position of a neuron relative to its mouse-mates is $-0.01$ to 0.11 over 1,259 cells. The decoder recovers where an animal's imaged patch looks, not which neuron within it looks where. $\Rsq$ on absolute coordinates then depends mostly on whether the between-mouse spread of a particular test set is reproduced at the right scale.

The within-mouse layout is present in the correlations, since same-mouse correlation falls with distance. It is not learned because each mouse contributes only 23 to 112 labelled cells. With 11,326 labels in one animal, MICrONS lets the same decoder reach neuron-level receptive fields at $\Rsq = 0.48$.

\begin{figure}[tbp]
  \centering
  \includegraphics[width=0.85\textwidth]{fig6_allen_rf_mouse_level.pdf}
  \caption{\textbf{Allen receptive fields: what crosses animals is each animal's screen position.} Predicted versus true mean receptive-field position of each held-out mouse (27 points over 3 splits; $x$ squares, $y$ dots), from decoders trained on single-animal (left) or mixed (right) populations. The per-mouse mean is recovered; the position of a neuron relative to its mouse-mates is not.}
  \label{fig:rf}
\end{figure}

\subsection{What moves the numbers}

More capacity helps where labels are plentiful and not where they are scarce (Fig.~\ref{fig:a1}). Twin receptive fields go from 0.34 to 0.48 to 0.51 $\Rsq$ between 2.2M, 17M and 57M parameters. Orientation error is flat from 2.2M upward on both MICrONS substrates ($20.9^\circ$, $20.2^\circ$, $19.9^\circ$ on the twin at 2.2M, 17M and 57M) and on Allen ($35.2^\circ$, $36.5^\circ$, $36.7^\circ$); only the 0.3M model is worse ($27.5^\circ$ and $27.8^\circ$ on MICrONS). Every model above 2M memorises its training neurons within 10 to 20 epochs and is early-stopped on the held-out selection slice (Fig.~\ref{fig:a3}). Population size beyond 256 to 512 neurons changes orientation error by at most $2^\circ$ in either direction, while twin receptive fields keep gaining slowly (Fig.~\ref{fig:a2}). Recomputing each training population's Gram from a random subset of the stimulus conditions, a label-free augmentation, gives no consistent gain on the Allen error ($35.2^\circ$, $33.5^\circ$, $37.7^\circ$ and $38.0^\circ$ with 100\%, 75\%, 50\% and 30\% of the conditions at 2.2M; $36.5^\circ$, $34.3^\circ$ and $35.6^\circ$ with 100\%, 85\% and 50\% at 17M; single runs, against a split-to-split spread of $3^\circ$) and is not used. Split-to-split variance on held-out mice exceeds seed variance (Fig.~\ref{fig:a4}).

\clearpage
\section{Discussion}

The correlation matrix of a few hundred cortical neurons, with the stimulus unknown and no neuron labelled, determines each neuron's preferred orientation to a mean error of $25^\circ$ in vivo and $20^\circ$ on the digital twin, against $45^\circ$ for an uninformed decoder and $41^\circ$ for the best constant guess. The readout is cardinal: near-cardinal neurons to $12^\circ$ to $19^\circ$, obliques near chance. It determines each neuron's receptive-field position to $\Rsq = 0.23$ in vivo and 0.48 on the digital twin. The symmetry analysis gives the mechanism: the relations fix the class structure up to the symmetries of the content, and a departure from that symmetry fixes the frame. That departure is a property of the correlations themselves: it is not the uneven number of neurons per orientation and not the particular stimulus sample. In each area the neurons preferring the over-represented orientation co-fluctuate most strongly, whether or not the sampled population is imbalanced. A direct ablation of that feature in the data is not available, but a model population with that co-fluctuation and nothing else, at the strength measured in cortex, recovers the absolute frame, while models with sharper or more numerous cardinal neurons do not.

Across animals the two contents behave differently. Orientation is carried by the relations between tuning profiles, which exist between any two neurons that saw the same conditions, in the same brain or not. From the within-circuit correlations of a single animal it is readable only weakly, at $31^\circ$ to $33^\circ$ on MICrONS, and on Allen not at all, where a single animal's 140 training cells are too few to learn from. Receptive-field position is carried at the level of an animal's patch, and the neuron-level layout needs more labels per animal than these recordings provide. Both statements concern what correlations under a shared stimulus contain, and both come from a protocol in which the test neurons, and in the held-out regimes the test animals, never touch training or model selection.

Limits. MICrONS is one animal, and its training and test halves share recording sessions. The Allen held-out-animal cells are the cleaner generalisation test, but their variance across mouse splits is large, and they exist only with grating relations, a designed 40-condition set on which two tuning profiles can be compared, not with the naturalistic relations of the MICrONS result. The readout is cardinal: the obliques are not resolved on either dataset, so the mean error understates the precision at the cardinals and overstates it elsewhere. The mechanism is established by exclusion, by cross-area transfer and by a model, not by removing the candidate anisotropy from cortical data, since the direct ablation we tried injected the labels it was meant to remove (Appendix~\ref{app:controls}). The model result is seed-sensitive near the cortical strength. Orientation labels are noisy in both datasets (in vivo and twin labels agree on 70\% of MICrONS neurons at 8 classes; static and drifting gratings on 45 to 67\% of Allen cells), and the Allen labels are quantised to $30^\circ$, which caps every labelled and label-free number alike. The twin substrate is a model fitted to the same recordings, so the in-vivo results carry the claim about cortex and the twin results show what cleaner correlations would give. The receptive-field labels come from the twin for both substrates, so the twin receptive-field result compares a model's correlations with the same model's receptive fields, and the in-vivo result scores cortical correlations against model-estimated positions. The orientation labels are in-vivo measurements throughout.

The symmetry result is an instance of a general expectation: a representation that is equivariant to a group has a group-fixed part, identical across individuals up to relabelling, and an individual remainder \citep{oizumi2026}. The matching-over-relabellings test is the discrete form of unsupervised alignment of similarity structures across individuals by optimal transport \citep{kawakita2024}. This paper adds the measurement in cortex, at the level of single neurons, and the finding that the absolute readout is possible and rests on a remainder that belongs to the correlations, not to the sample of neurons or stimuli.

\paragraph{Code and data.} TODO: repository URL. All data are public releases.

\bibliographystyle{unsrtnat}
\bibliography{refs}

\clearpage
\appendix
\renewcommand{\thefigure}{A\arabic{figure}}
\setcounter{figure}{0}
\section{Scaling and controls}
Figures~\ref{fig:a1} to \ref{fig:a4} report the sweeps behind Section~3.5: decoder size, population size, training dynamics and split variance. Two training-side variants were tried and dropped: recomputing each training population's Gram from a random subset of the stimulus conditions, and zeroing random Gram entries. Neither changes the angular error by more than the split-to-split spread (Section~3.5).

\section{Controls}
\label{app:controls}
Four controls on the MICrONS standard configuration (17M, 512-neuron populations, digital twin unless stated; Table~\ref{tab:controls}).

\paragraph{What the decoder reads.} With the three row statistics as the only input per neuron and no relational term, the error is $37.7^\circ$, raw and after alignment alike, so the scalars fix a weak frame on their own. They are not a selectivity confound: their correlations with the distance to the nearest cardinal axis are below 0.12 in magnitude and with gOSI below 0.16, and gOSI is 0.39 for cardinal and 0.37 for oblique neurons. With the statistics plus the Gram entering only as an attention bias, which makes the decoder equivariant to the order of the neurons, the error is $19.9^\circ$ over three seeds, the same as the standard decoder ($20.2^\circ$), whose input layer also projects the raw Gram row in sample order. The relations carry the readout, and the order-dependent input layer contributes nothing.

\paragraph{Reference decoders that see the activity.} To measure what the relations discard, the same transformer was trained with each neuron's token built from its response vector instead of its Gram row, so that it sees the stimulus-aligned activity the Gram is computed from (Table~\ref{tab:ref}). It does better than the relational decoder on both contents and both substrates. Adding the Gram to it as an attention bias changes nothing. A linear readout of each neuron's own response vector, with no population context, does as well as the activity transformer, so the activity decoder is a per-neuron tuning readout whose only advantage over the relational decoder is the alignment of the responses to the stimulus. On orientation the relations retain about 85\% of the gain over chance that the aligned activity gives, and on receptive fields about 60\%; the twin receptive-field gap is inflated by the twin's labels and responses coming from the same model. In exchange the relational decoder does not depend on the stimulus alignment. With test responses from movie frames the decoder never trained on, the activity decoder has no valid input, since its coordinates now index different frames, and it falls to or below chance, while the relational decoder scores $33.8^\circ$ in vivo and $20.7^\circ$ on the twin. This says what a stimulus-aligned readout can be applied to, not that it lost a contest. The activity decoder transfers across animals under a shared stimulus protocol (Section~3.3) but not across stimuli. The relations do both, because a correlation does not depend on which frames it was computed from, only on both neurons having seen the same ones. We also asked whether the pairwise Gram is all the relations offer. A decoder given the response vectors with the stimulus bins shuffled, the same shuffle for every neuron of a population, is blind to the stimulus but free to compute any relation among the neurons, pairwise or higher. On orientation it is no better than the Gram decoder over three seeds and, on the twin, unstable across them ($22^\circ$ to $37^\circ$); on receptive fields it is far below the Gram decoder. The shuffled vector without population context is near chance, so what the shuffled decoder does read comes from relations it computes itself. We find no evidence that a decoder free to compute richer statistics extracts more than the pairwise Gram, with the caveat that it has to discover the relations under a fresh permutation for every population, a harder problem than being handed them. The comparison does show that giving the decoder the relations explicitly is better than leaving it to find them.

\begin{table}[htbp]
  \centering\footnotesize
  \caption{Reference decoders that see the activity (each neuron's stimulus-aligned response vector), 17M unless linear, 512-neuron populations. Orientation as mean angular error (chance $45^\circ$), receptive field as $\Rsq$; mean $\pm$ s.d.\ over three seeds where given, otherwise seed 0.}
  \label{tab:ref}
  \begin{tabular}{lcccc}
    \toprule
    decoder input & ori.\ in vivo & ori.\ twin & RF in vivo & RF twin \\
    \midrule
    Gram row (this paper) & $25.4^\circ \pm 0.5^\circ$ & $20.2^\circ \pm 0.4^\circ$ & $0.23 \pm 0.01$ & $0.48 \pm 0.01$ \\
    activity, transformer & $21.5^\circ \pm 0.4^\circ$ & $15.0^\circ \pm 0.1^\circ$ & 0.37 & 0.80 \\
    activity + Gram as attention bias & $22.2^\circ$ & $15.0^\circ$ & -- & -- \\
    activity, linear, no population context & $20.2^\circ \pm 0.2^\circ$ & $15.8^\circ \pm 0.1^\circ$ & $0.34 \pm 0.01$ & $0.68 \pm 0.01$ \\
    activity, transformer, disjoint stimulus bins & $49.5^\circ$ & $41.0^\circ$ & -- & -- \\
    activity, bins shuffled per population, transformer & $26.3^\circ \pm 0.6^\circ$ & $27.8^\circ \pm 6.5^\circ$ & 0.04 & 0.08 \\
    activity, bins shuffled, linear, no population context & $40.3^\circ$ & $40.2^\circ$ & -- & -- \\
    \bottomrule
  \end{tabular}
\end{table}

\paragraph{Cross-stimulus test.} Training and selection Grams built from one random half of the stimulus bins and test Grams from the other half give $20.7^\circ$ on the twin (2,500 bins per half, raw responses rather than principal components; three seeds, 20.3 to 21.2) and $33.8^\circ$ in vivo (60 bins per half; 33.4 to 34.2), against $20.2^\circ$ and $25.4^\circ$ with shared bins. A control with the same 60-bin half on both sides gives $25.3^\circ$, so the in-vivo increase comes from the stimulus sample, not the smaller Gram: a correlation matrix over 60 movie bins partly reflects which frames it was built from, and one over 2,500 bins does not. The relational structure the decoder reads is stable across stimuli once the stimulus sample is large.

\paragraph{Scan and area.} The MICrONS neurons come from 13 scans in four areas, same-scan pairs correlate more strongly than cross-scan pairs, and the scans look at different parts of the screen, so a decoder that read scan identity from the Gram could predict each scan's mean receptive-field position and score a positive $\Rsq$ without knowing any neuron's position. The Allen decoder does exactly that (Section~3.4), so we checked MICrONS the same way, centring predictions and labels per scan and per area before scoring. The receptive-field result holds: on a fourth seed of the Table~\ref{tab:microns} configuration, run with saved predictions, $\Rsq$ is 0.50 absolute and 0.50 within scan on the twin, 0.25 and 0.25 in vivo, and 0.47 and 0.22 within area. Only 2\% of the label variance lies between scans and 9\% between areas, so a scan-identity readout could reach at most 0.02. For orientation, a predictor that assigns every neuron its scan's or area's mean orientation scores $39^\circ$, and removing a separate rotation per scan or area changes the decoder's error by less than $0.2^\circ$. Both MICrONS results are neuron-level readouts within a scan.

\paragraph{Class balance and priors.} Preferred orientations are uneven in MICrONS (19\% of labelled neurons within $11^\circ$ of $0^\circ$, 25\% of $90^\circ$, 6 to 10\% in each oblique bin), so a decoder could gain from the distribution alone. It gains little: the best constant prediction has an error of $41.3^\circ$, a predictor that assigns every neuron the circular mean of its scan $39.0^\circ$ and of its area $38.7^\circ$, against $20.3^\circ$ for the decoder on the twin and $24.6^\circ$ in vivo (neuron split 2). Training and testing on populations drawn with equal numbers of neurons from each of the eight classes (64 per class) gives $22.0^\circ$ on the twin and $26.6^\circ$ in vivo (two seeds: 22.4, 21.5; 26.5, 26.7), raw and after global alignment alike, so the absolute frame does not depend on the class imbalance. The per-area errors are within $3^\circ$ of each other. On Allen the labels are nearly flat (14 to 20\% per orientation), the best constant is $43.9^\circ$, and the balanced error is $36.4^\circ$ against $35.7^\circ$ raw on the same saved predictions (held-out-animal split 0). The per-bin errors show, on both datasets, that the readout is cardinal: neurons preferring $0^\circ$ or $90^\circ$ are decoded to $11^\circ$ to $12^\circ$ on the twin and $25^\circ$ on Allen, oblique ones to $37^\circ$ to $45^\circ$, and 63 to 72\% of all predictions fall on the two cardinal axes (Fig.~\ref{fig:bins}). The relations say how close a neuron is to a cardinal axis, the anisotropy of Section~3.2 seen from the decoder's side, and much less about which oblique it prefers.

\begin{table}[htbp]
  \centering\footnotesize
  \caption{Controls, mean angular error on held-out neurons (chance $45^\circ$), 17M decoder on 512-neuron populations unless stated; single runs unless a seed count is given. Where two numbers are given, the second is the error after the best global rotation or reflection (rows marked \emph{aligned}) or averaged over eight equally weighted true-orientation bins (rows marked \emph{balanced}).}
  \label{tab:controls}
  \begin{tabular}{llc}
    \toprule
    control & substrate & error \\
    \midrule
    standard decoder & twin & $20.2^\circ$ \\
    row statistics only, no relational term & twin & $37.7^\circ$ \\
    row statistics + Gram as attention bias (3 seeds) & twin & $19.9^\circ \pm 0.3^\circ$ \\
    disjoint stimulus bins, training vs test (3 seeds) & twin & $20.7^\circ \pm 0.4^\circ$ \\
    standard decoder & in vivo & $25.4^\circ$ \\
    disjoint stimulus bins (60 per side; 3 seeds) & in vivo & $33.8^\circ \pm 0.3^\circ$ \\
    same 60 bins on both sides & in vivo & $25.3^\circ$ \\
    neuron split 1 / 2 & in vivo & $25.4^\circ$ / $24.6^\circ$ \\
    standard decoder, raw / balanced & twin & $20.3^\circ$ / $24.5^\circ$ \\
    orientation-balanced loss, raw / balanced & twin & $21.4^\circ$ / $25.3^\circ$ \\
    standard decoder, raw / balanced & in vivo & $24.6^\circ$ / $27.9^\circ$ \\
    orientation-balanced loss, raw / balanced & in vivo & $25.8^\circ$ / $29.1^\circ$ \\
    class-balanced populations (64 per class, 2 seeds), raw / aligned & twin & $22.0^\circ$ / $21.9^\circ$ \\
    class-balanced populations, raw / aligned, 2 seeds & in vivo & $26.6^\circ$ / $26.4^\circ$ \\
    trained on V1, scored on V1 (128 neurons), raw / aligned & twin & $22.2^\circ$ / $22.2^\circ$ \\
    trained on V1, scored on RL (512 neurons), raw / aligned & twin & $26.2^\circ$ / $24.6^\circ$ \\
    trained on V1, scored on AL (128 neurons), raw / aligned & twin & $58.2^\circ$ / $31.6^\circ$ \\
    within-scan populations (128 neurons) & twin / in vivo & $30.7^\circ$ / $32.6^\circ$ \\
    one training scan (140 neurons), populations of 128 & twin / in vivo & $38^\circ$--$44^\circ$ / $37^\circ$--$39^\circ$ \\
    mixed-scan populations (128 neurons, V1) & twin & $22.2^\circ$ \\
    \midrule
    Allen, mixed populations of 128, training pool 280 / 560 / 6,706 & 2.2M & $40.5^\circ$ / $40.7^\circ$ / $36.0^\circ$ \\
    \bottomrule
  \end{tabular}
\end{table}

\paragraph{A subtraction ablation that leaks labels.} To remove the class-level anisotropy from the input, we estimated the class-Gram residual (8 or 24 classes) on the labelled training half and subtracted the class-pair offset from every Gram entry, using each neuron's true class, at training and test. The twin error was unchanged ($20.7^\circ$ raw and aligned at 8 classes) and improved at 24 classes ($17.7^\circ$); the control that subtracts the whole class structure, which should destroy decoding, improved it further ($17.2^\circ$ raw, $15.6^\circ$ aligned). An offset keyed to the true class is a label written into the Gram, through its jumps at class boundaries and its estimation noise, and the decoder reads it. No result of this ablation is used.

\paragraph{Synthetic populations.} 6,000 model neurons with von Mises tuning ($\kappa = 2$), preferred orientations either uniform or with half the neurons within about $12^\circ$ of a cardinal axis (count), tuning either isotropic or sharpened to $\kappa = 4$ at the cardinals (sharp), and a 120-frame stimulus that either covers orientation evenly or has half its frames near a cardinal axis (stim); responses with Gaussian noise; the 2.2M decoder on 512-neuron populations, epoch selected on the raw error as in every other run, five data and decoder seeds (Table~\ref{tab:syn}). With every switch off the class-Gram is circulant to within noise (neighbour-correlation range 0.01 against 0.04 to 0.14 with a switch on), and the decoder recovers the ordering but no frame. Sharper cardinal tuning and a cardinal excess of neurons fix the frame in at most one seed of five, alone, together, or at any strength ($\kappa$ up to 32; gain up to 5 times; two seeds each, all above $42^\circ$ raw). Stimulus anisotropy fixes it in most seeds. The anisotropy the cortical data show, stronger co-fluctuation among neurons that prefer the area's dominant orientation (Section~3.2), is modelled by a smooth random field over preferred orientation entering each neuron with a weight peaked at $90^\circ$, with uniform counts, isotropic tuning and an even stimulus. The anisotropy factor of Section~3.2, measured on the model over its five data seeds, is 1.09, 1.32 and 1.68 at the three strengths run. The middle one matches the twin's V1 (1.28) and fixes the frame in four of five seeds, the strongest in both seeds, the weakest, below any cortical value, in one of two. The spread across seeds is large in every cell where the frame is not fixed, because selection on the raw error then has no signal to select on. The cell means are therefore reported with their range. The frame-fixed counts are unchanged with the raw-error threshold at $25^\circ$ or $35^\circ$, except that the count cell drops to 0 of 5 at $25^\circ$. One run (all three switches, seed 0) never left its initial loss and recovered neither ordering nor frame; it is included in the mean.

\begin{table}[htbp]
  \centering\small
  \caption{Synthetic populations, five seeds: raw error / error after the best global rotation or reflection, mean over seeds with the range in brackets, and the number of seeds in which the frame was fixed (raw error below $30^\circ$ and within $5^\circ$ of the aligned error). Bottom block: shared variability among like-tuned neurons entering with a weight peaked at $90^\circ$; the factor is the anisotropy factor of Section~3.2 measured on the model. A gap between raw and aligned error means structure without frame.}
  \label{tab:syn}
  \begin{tabular}{lccc}
    \toprule
    anisotropy switched on & raw & aligned & frame fixed \\
    \midrule
    none (circulant) & $53.7^\circ$ [44--63] & $30.0^\circ$ [23--42] & 0 / 5 \\
    stimulus & $23.5^\circ$ [4--54] & $13.9^\circ$ [4--19] & 3 / 5 \\
    sharper cardinal tuning & $51.4^\circ$ [39--65] & $28.6^\circ$ [15--42] & 0 / 5 \\
    cardinal excess of neurons (count) & $37.8^\circ$ [27--45] & $35.1^\circ$ [27--42] & 1 / 5 \\
    sharper tuning + stimulus & $12.2^\circ$ [3--38] & $6.3^\circ$ [3--10] & 4 / 5 \\
    count + stimulus & $25.0^\circ$ [4--47] & $20.7^\circ$ [4--43] & 3 / 5 \\
    count + sharper tuning & $42.8^\circ$ [29--58] & $35.2^\circ$ [27--41] & 1 / 5 \\
    all three & $18.4^\circ$ [3--46] & $17.9^\circ$ [3--44] & 4 / 5 \\
    \midrule
    like-tuned co-fluctuation, factor 1.32 (twin V1: 1.28) & $8.9^\circ$ [4--26] & $7.3^\circ$ [3--19] & 4 / 5 \\
    \quad factor 1.09 (two seeds) & $31.5^\circ$ [15--48] & $24.2^\circ$ [15--33] & 1 / 2 \\
    \quad factor 1.68 (two seeds) & $2.6^\circ$ [2--3] & $2.5^\circ$ [2--3] & 2 / 2 \\
    \quad factor 1.32 + stimulus (two seeds) & $5.2^\circ$ [4--6] & $5.2^\circ$ [4--6] & 2 / 2 \\
    \bottomrule
  \end{tabular}
\end{table}

\section{Classification results}
Earlier runs decoded orientation as a class (8 classes on MICrONS, 6 on Allen) with the same decoder and protocol. The per-neuron accuracies are in Table~\ref{tab:class}. They lead to the same conclusions, but a class score counts a near miss as a miss, so the angular error in the main text is the better summary of a continuous label.

\begin{table}[htbp]
  \centering\small
  \caption{Orientation as classification, 17M decoder. MICrONS: 8 classes, majority class 0.255, mean $\pm$ s.d.\ over 3 seeds. Allen: 6 classes, majority class 0.19 to 0.20, three splits per cell (17M decoder with the condition-subset augmentation, 256-neuron populations, not the runs of Table~\ref{tab:allen}; the \code{pooledcross} split-0 value is the mean of 3 seeds).}
  \label{tab:class}
  \begin{tabular}{llc}
    \toprule
    dataset & cell & accuracy \\
    \midrule
    MICrONS & in vivo & $0.380 \pm 0.006$ \\
    MICrONS & digital twin & $0.420 \pm 0.003$ \\
    Allen & training animals, single-animal populations & 0.210 / 0.216 / 0.217 \\
    Allen & training animals, mixed populations & 0.281 / 0.283 / 0.283 \\
    Allen & held-out animals, single-animal populations & 0.199 / 0.197 / 0.187 \\
    Allen & held-out animals, mixed populations & 0.293 / 0.236 / 0.264 \\
    \bottomrule
  \end{tabular}
\end{table}

\section{Labelled reference models}
A second reference uses labels at test time: a ridge readout from each test neuron's correlations to the labelled training neurons, fitted on the training half, from the classification-era runs (orientation as 8 or 6 classes, accuracy; receptive fields as $\Rsq$). It is one linear labelled model, not a ceiling, and its orientation column is an accuracy, not comparable with the angular errors of the main text.

\begin{table}[htbp]
  \centering\small
  \caption{Ridge readout from a neuron's correlations to labelled training neurons (classification era; orientation accuracy, receptive field $\Rsq$).}
  \begin{tabular}{lcc}
    \toprule
    dataset, substrate & orientation & receptive field \\
    \midrule
    MICrONS, in vivo & 0.41 (8 classes) & 0.32 \\
    MICrONS, digital twin & 0.49 (8 classes) & 0.67 \\
    Allen, leave-one-mouse-out & 0.35 (6 classes) & 0.20 \\
    \bottomrule
  \end{tabular}
\end{table}

\section{Receptive-field decomposition per split}
Per held-out-mouse split, for the two Allen regimes of Fig.~\ref{fig:rf}: $\Rsq$ on absolute screen coordinates, the correlation between predicted and true per-mouse mean position over the 9 test mice ($x$, $y$), and the correlation between predicted and true position relative to the mouse mean over the cells ($x$, $y$).

\begin{table}[htbp]
  \centering\small
  \caption{Allen receptive fields per split.}
  \begin{tabular}{llccc}
    \toprule
    populations & split & absolute $\Rsq$ & per-mouse mean $r$ ($x$, $y$) & within-mouse $r$ ($x$, $y$) \\
    \midrule
    single-animal & 0 & 0.29 & 0.47, 0.74 & 0.07, 0.11 \\
    single-animal & 1 & $-0.02$ & 0.48, 0.40 & 0.06, 0.08 \\
    single-animal & 2 & $-0.45$ & 0.19, 0.25 & 0.11, $-0.01$ \\
    mixed & 0 & 0.22 & 0.70, 0.81 & 0.00, 0.20 \\
    mixed & 1 & 0.08 & 0.23, 0.80 & 0.02, 0.08 \\
    mixed & 2 & $-0.25$ & 0.46, 0.75 & $-0.06$, 0.04 \\
    \bottomrule
  \end{tabular}
\end{table}


\clearpage
\begin{figure}[p]
  \centering
  \includegraphics[width=\textwidth]{figA1_scaling_params.pdf}
  \caption{\textbf{Scaling with decoder size.} Test metric versus number of decoder parameters at fixed population size (512 neurons MICrONS, 256 Allen). Error bars: s.d.\ over 3 seeds where run. On MICrONS receptive fields the gain continues to 57M; the orientation error is flat from 2.2M upward on both substrates and on Allen.}
  \label{fig:a1}
\end{figure}

\begin{figure}[p]
  \centering
  \includegraphics[width=0.85\textwidth]{figA2_population_size.pdf}
  \caption{\textbf{Population size.} Test metric versus the number of neurons in each sampled population. Beyond 256 to 512 neurons the orientation error moves by at most $2^\circ$, up or down; twin receptive fields keep gaining slowly, Allen receptive fields lose.}
  \label{fig:a2}
\end{figure}

\begin{figure}[p]
  \centering
  \includegraphics[width=\textwidth]{figA3_training_dynamics.pdf}
  \caption{\textbf{Training dynamics and epoch selection.} Test metric (solid), selection-set metric (dotted) and training loss (red, right axis) per epoch. The selection set, never the test set, picks the epoch (grey line). The training loss collapses within 10 to 20 epochs on every large run.}
  \label{fig:a3}
\end{figure}

\begin{figure}[p]
  \centering
  \includegraphics[width=\textwidth]{figA4_allen_splits_classlevel.pdf}
  \caption{\textbf{Allen split variance and the class-level test.} Left: \code{pooledcross} orientation error on the same three held-out-mouse splits for the previous standard configuration (17M, augmented) and the final one (2.2M, plain); split 0 carries three seeds. Right: the class-level test with relations from natural movies versus drifting gratings; only the grating relations move every test off chance.}
  \label{fig:a4}
\end{figure}


\end{document}
